\documentclass[10pt,a4paper]{amsart}
\usepackage{amsmath,amssymb,amsthm,graphicx}
\usepackage{verbatim}
\newenvironment{reference}{\comment}{\endcomment}
%\newenvironment{reference}{}{}
\newenvironment{slogan}{\comment}{\endcomment}
\newenvironment{history}{\comment}{\endcomment}

% For commutative diagrams we use Xy-pic
\usepackage[all]{xy}

% We use 2cell for 2-commutative diagrams.
\xyoption{2cell}
% Theorem environments.
%
\theoremstyle{plain}
\newtheorem{theorem}[subsection]{Theorem}
\newtheorem{proposition}[subsection]{Proposition}
\newtheorem{lemma}[subsection]{Lemma}

\theoremstyle{definition}
\newtheorem{definition}[subsection]{Definition}
\newtheorem{example}[subsection]{Example}
\newtheorem{exercise}[subsection]{Exercise}
\newtheorem{situation}[subsection]{Situation}

\theoremstyle{remark}
\newtheorem{remark}[subsection]{Remark}
\newtheorem{remarks}[subsection]{Remarks}

\numberwithin{equation}{subsection}

% Macros
%
\def\lim{\mathop{\mathrm{lim}}\nolimits}
\def\colim{\mathop{\mathrm{colim}}\nolimits}
\def\Spec{\mathop{\mathrm{Spec}}}
\def\Hom{\mathop{\mathrm{Hom}}\nolimits}
\def\Ext{\mathop{\mathrm{Ext}}\nolimits}
\def\SheafHom{\mathop{\mathcal{H}\!\mathit{om}}\nolimits}
\def\SheafExt{\mathop{\mathcal{E}\!\mathit{xt}}\nolimits}
\def\Sch{\mathit{Sch}}
\def\Mor{\mathop{\mathrm{Mor}}\nolimits}
\def\Ob{\mathop{\mathrm{Ob}}\nolimits}
\def\Sh{\mathop{\mathit{Sh}}\nolimits}
\def\NL{\mathop{N\!L}\nolimits}
\def\CH{\mathop{\mathrm{CH}}\nolimits}
\def\proetale{{pro\text{-}\acute{e}tale}}
\def\etale{{\acute{e}tale}}
\def\QCoh{\mathit{QCoh}}
\def\Ker{\mathop{\mathrm{Ker}}}
\def\Im{\mathop{\mathrm{Im}}}
\def\Coker{\mathop{\mathrm{Coker}}}
\def\Coim{\mathop{\mathrm{Coim}}}

% Boxtimes
%
\DeclareMathSymbol{\boxtimes}{\mathbin}{AMSa}{"02}

%
% Macros for moduli stacks/spaces
%
\def\QCohstack{\mathcal{QC}\!\mathit{oh}}
\def\Cohstack{\mathcal{C}\!\mathit{oh}}
\def\Spacesstack{\mathcal{S}\!\mathit{paces}}
\def\Quotfunctor{\mathrm{Quot}}
\def\Hilbfunctor{\mathrm{Hilb}}
\def\Curvesstack{\mathcal{C}\!\mathit{urves}}
\def\Polarizedstack{\mathcal{P}\!\mathit{olarized}}
\def\Complexesstack{\mathcal{C}\!\mathit{omplexes}}
% \Pic is the operator that assigns to X its picard group, usage \Pic(X)
% \Picardstack_{X/B} denotes the Picard stack of X over B
% \Picardfunctor_{X/B} denotes the Picard functor of X over B
\def\Pic{\mathop{\mathrm{Pic}}\nolimits}
\def\Picardstack{\mathcal{P}\!\mathit{ic}}
\def\Picardfunctor{\mathrm{Pic}}
\def\Deformationcategory{\mathcal{D}\!\mathit{ef}}
\usepackage{cleveref}
\crefname{originallemma}{original lemma}{original lemmas}
\Crefname{originallemma}{Original lemma}{Original lemmas}
\makeatletter
\newlabel{algebra-lemma-completion-tensor}{{00MA}{source}{}{sourcepointer-1}{}}
\newlabel{algebra-lemma-completion-tensor@cref}{{[originallemma][0][]00MA}{[1][0][]source}{}{}{}}
\makeatother
\crefname{originallemma}{original lemma}{original lemmas}
\Crefname{originallemma}{Original lemma}{Original lemmas}
\makeatletter
\newlabel{algebra-lemma-flat-base-change}{{00HI}{source}{}{sourcepointer-2}{}}
\newlabel{algebra-lemma-flat-base-change@cref}{{[originallemma][0][]00HI}{[1][0][]source}{}{}{}}
\makeatother
\crefname{originalreference}{original reference}{original references}
\Crefname{originalreference}{Original reference}{Original references}
\makeatletter
\newlabel{lemma-rig-closed-point-after-localization}{{0GGG}{source}{}{sourcepointer-3}{}}
\newlabel{lemma-rig-closed-point-after-localization@cref}{{[originalreference][0][]0GGG}{[1][0][]source}{}{}{}}
\makeatother
\crefname{originalreference}{original reference}{original references}
\Crefname{originalreference}{Original reference}{Original references}
\makeatletter
\newlabel{lemma-rig-closed-point-relative}{{0GGB}{source}{}{sourcepointer-4}{}}
\newlabel{lemma-rig-closed-point-relative@cref}{{[originalreference][0][]0GGB}{[1][0][]source}{}{}{}}
\makeatother
\crefname{originalreference}{original reference}{original references}
\Crefname{originalreference}{Original reference}{Original references}
\makeatletter
\newlabel{lemma-rig-closed-point-variables}{{0GGH}{source}{}{sourcepointer-5}{}}
\newlabel{lemma-rig-closed-point-variables@cref}{{[originalreference][0][]0GGH}{[1][0][]source}{}{}{}}
\makeatother
\crefname{originallemma}{original lemma}{original lemmas}
\Crefname{originallemma}{Original lemma}{Original lemmas}
\makeatletter
\newlabel{more-algebra-lemma-flat-completion}{{0C4G}{source}{}{sourcepointer-6}{}}
\newlabel{more-algebra-lemma-flat-completion@cref}{{[originallemma][0][]0C4G}{[1][0][]source}{}{}{}}
\makeatother
\title{Rig-flatness is stable under adic topologically finite type base change}
\author{The Stacks Project Authors}
\date{stacks-project:restricted.tex:local-snapshot-20261012; GNU FDL 1.2 or later}
\begin{document}
\maketitle
\noindent\textit{Editorial note.} All delivered mathematical text is the author's own, copied from the pinned Stacks Project chapter file restricted.tex; no statement or proof step was rewritten, reordered or removed.\par\medskip
\noindent\textit{Editorial note.} The counted proof is the author's complete argument: the author also flags the identity as something used below; it is the compressed standard fact, not a verification specific to the claim\par\medskip
\noindent\textit{Editorial note.} The label lemma-rig-flat-base-change names the counted statement, and the proof environment immediately following it in the source is delivered verbatim as its proof.\par\medskip
\noindent\textit{Editorial note.} The counted statement is carried by the author's own lemma environment in the chapter source.\par\medskip
\noindent\textit{Allowed standard prerequisites (editorial).}\begin{itemize}
\item Lemma lemma-rig-closed-point-relative (delivered, L48108): contraction behaviour of rig-closed primes through adic topologically finite type maps
\item Lemma lemma-rig-closed-point-variables (in-source): at rig-closed primes the completions of the restricted and polynomial power series local rings are isomorphic
\item More on Algebra, Lemma more-algebra-lemma-flat-completion: flatness transfers through isomorphisms on completions
\item Algebra, Lemma algebra-lemma-completion-tensor and algebra-lemma-flat-base-change
\item standard formal-space identity A\_\{f\}\{x\} = A\{x\}\_\{f: localization commutes with restricted power series
\end{itemize}\medskip
\ifdefined\hypertarget\hypertarget{sourcepointer-1}{}\fi
\noindent\textit{Source reference.} Original lemma 00MA: sources/stacks-chow/tags.txt, source lines 816--816. Prerequisite pointer invoked at 3819; it names a standard fact of a sibling chapter that the counted argument uses, and every step of the counted proof built on it is retained verbatim in the delivered body.\par\medskip
\ifdefined\hypertarget\hypertarget{sourcepointer-2}{}\fi
\noindent\textit{Source reference.} Original lemma 00HI: sources/stacks-chow/tags.txt, source lines 644--644. Prerequisite pointer invoked at 3821; it names a standard fact of a sibling chapter that the counted argument uses, and every step of the counted proof built on it is retained verbatim in the delivered body.\par\medskip
\ifdefined\hypertarget\hypertarget{sourcepointer-3}{}\fi
\noindent\textit{Source reference.} Original reference 0GGG: sources/stacks-restricted/pinned-original-source.tex, source lines 3488--3488. Prerequisite pointer invoked at 3842; it names a standard fact of the same chapter that the counted argument uses, and every step of the counted proof built on it is retained verbatim in the delivered body.\par\medskip
\ifdefined\hypertarget\hypertarget{sourcepointer-4}{}\fi
\noindent\textit{Source reference.} Original reference 0GGB: sources/stacks-restricted/pinned-original-source.tex, source lines 3333--3333. Prerequisite pointer invoked at 3868; it names a standard fact of the same chapter that the counted argument uses, and every step of the counted proof built on it is retained verbatim in the delivered body.\par\medskip
\ifdefined\hypertarget\hypertarget{sourcepointer-5}{}\fi
\noindent\textit{Source reference.} Original reference 0GGH: sources/stacks-restricted/pinned-original-source.tex, source lines 3538--3538. Prerequisite pointer invoked at 3881; it names a standard fact of the same chapter that the counted argument uses, and every step of the counted proof built on it is retained verbatim in the delivered body.\par\medskip
\ifdefined\hypertarget\hypertarget{sourcepointer-6}{}\fi
\noindent\textit{Source reference.} Original lemma 0C4G: sources/stacks-chow/tags.txt, source lines 15226--15226. Prerequisite pointer invoked at 3888; it names a standard fact of a sibling chapter that the counted argument uses, and every step of the counted proof built on it is retained verbatim in the delivered body.\par\medskip
\begin{lemma}
\label{lemma-rig-flat-base-change}
Let $\varphi : A \to B$ and $A \to C$ be arrows of $\textit{WAdm}^{Noeth}$.
Assume $\varphi$ is rig-flat and $A \to C$ adic and topologically
of finite type. Then $C \to B \widehat{\otimes}_A C$ is rig-flat.
\end{lemma}

\begin{proof}
Assume $\varphi$ is rig-flat. Let $f \in C$ be an element.
We have to show that
$C_{\{f\}} \to B \widehat{\otimes}_A C_{\{f\}}$
is naively rig-flat. Since we can replace $C$ by $C_{\{f\}}$
we it suffices to show that $C \to B \widehat{\otimes}_A C$
is naively rig-flat.

\medskip\noindent
If $A \to C$ is surjective or more generally if $C$ is finite as
an $A$-module, then $B \otimes_A C = B \widehat{\otimes}_A C$
as a finite module over a complete Noetherian ring is complete, see
Algebra, Lemma \ref{algebra-lemma-completion-tensor}.
By the usual base change for flatness
(Algebra, Lemma \ref{algebra-lemma-flat-base-change})
we see that naive rig-flatness of $\varphi$ implies naive rig-flatness
for $C \to B \times_A C$ in this case.

\medskip\noindent
In the general case, we can factor $A \to C$ as
$A \to A\{x_1, \ldots, x_n\} \to C$
where $A\{x_1, \ldots, x_n\}$ is the restricted power series ring
and $A\{x_1, \ldots, x_n\} \to C$ is surjective. Thus it
suffices to show $C \to B \widehat{\otimes}_A B$ is naively
rig-flat in case $C = A\{x_1, \ldots, x_n\}$.
Since $A\{x_1, \ldots, x_n\} = A\{x_1, \ldots, x_{n - 1}\}\{x_n\}$
by induction on $n$ we reduce to the case discussed in the next
paragraph.

\medskip\noindent
Here $C = A\{x\}$. Note that $B \widehat{\otimes}_A C = B\{x\}$.
We have to show that $A\{x\} \to B\{x\}$ is naively rig-flat.
Let $\mathfrak q \subset B\{x\}$ be a rig-closed prime ideal.
We have to show that $B\{x\}_{\mathfrak q}$ is flat over $A\{x\}$.
Set $\mathfrak p = A \cap \mathfrak q$.
By Lemma \ref{lemma-rig-closed-point-after-localization}
we can find an $f \in A$ such that
$f$ maps to a unit in $B\{x\}/\mathfrak q$ and such
that the prime ideal $\mathfrak p'$ in $A_{\{f\}}$ induced is rig-closed.
Below we will use that $A_{\{f\}}\{x\} = A\{x\}_{\{f\}}$ and similarly
for $B$; details omitted. Consider the diagram
$$
\xymatrix{
(B\{x\})_{\mathfrak q} \ar[r] &
(B_{\{f\}}\{x\})_{\mathfrak q'} \\
A\{x\} \ar[r] \ar[u] &
A_{\{f\}}\{x\} \ar[u]
}
$$
We want to show that the left vertical arrow is flat.
The top horizontal arrow is faithfully flat as it is a local
homomorphism of local rings and flat as $B_{\{f\}}\{x\}$
is the completion of a localization of the Noetherian ring
$B\{x\}$. Similarly the bottom horizontal arrow is flat.
Hence it suffices to prove that the right vertical arrow is flat.
This reduces us to the case discussed in the next paragraph.

\medskip\noindent
Here $C = A\{x\}$, we have a rig-closed prime ideal
$\mathfrak q \subset B\{x\}$ such that
$\mathfrak p = A \cap \mathfrak q$ is rig-closed as well.
This implies, via Lemma \ref{lemma-rig-closed-point-relative},
that the intermediate primes $B \cap \mathfrak q$ and
$A\{x\} \cap \mathfrak q$ are rig-closed as well.
Consider the diagram
$$
\xymatrix{
(B[x])_{B[x] \cap \mathfrak q} \ar[r] &
(B\{x\})_{\mathfrak q} \\
(A[x])_{A[x] \cap \mathfrak q} \ar[r] \ar[u] &
(A\{x\})_{A\{x\} \cap \mathfrak q} \ar[u]
}
$$
of local homomorphisms of Noetherian local rings.
By Lemma \ref{lemma-rig-closed-point-variables}
the horizontal arrows define isomorphisms
on completions. We already know that the left
vertical arrow is flat (as $A \to B$ is naively rig-flat
and hence $A[x] \to B[x]$ is flat away from the
closed locus defined by an ideal of definition).
Hence we finally conclude by
More on Algebra, Lemma \ref{more-algebra-lemma-flat-completion}.
\end{proof}

\end{document}
